\documentclass{report}
\usepackage{amsmath,amssymb}
\setlength{\parindent}{0mm}
\setlength{\parskip}{1em}
\begin{document}
\begin{center}
\rule{6in}{1pt} \
{\large Christian Ketelsen \\
{\bf An Element-Free Adaptive Algebraic Multigrid Method for Markov Chain Monte Carlo Simulations}}

Center for Applied Scientific Computing \\ Lawrence Livermore National Laboratory \\ Box 808 \\ L-560 \\ Livermore \\ CA 94551-0808
\\
{\tt ketelsen1@llnl.gov}\\
Panayot S.  Vassilevski\end{center}

We are interested in solving a $2^{\text{nd}}$ -order elliptic equation
where there is some uncertainty in the
background field, $a$:

\begin{equation}\label{PDE}
\mathcal{L}(a(x,\omega),\;u(x,\omega)) =f,
\end{equation}

subject to some boundary conditions. Monte Carlo methods have, for some
time, been the standard solution
method for such problems. In this setting, background field $a(x,\omega)$
is assumed to come from a prescribed
probability distribution. A great many samples of $a$ are drawn and the
forward problem is solved to obtain an
ensemble of solutions $\{u(x,\;\omega_i)\}^N_{i=1}$. Moments of the
ensemble can then be computed to quantify uncertainty in the solution.

Often times, background field $a$ belongs to a probability distribution
that is difficult to sample from.
Such is the case in simulations of subsurface flow and quantum field
theory (QFT). In such cases, sampling
is achieved using the computationally expensive Markov Chain Monte Carlo
(MCMC) method. In MCMC,
samples are obtained via an accept-reject process. Given the current
realization $a(x, \omega_k)$, proposed realization $a(x, \omega{'}_k)$ is
generated via some simple transition rule (perturbing the value in each
grid cell with a normal random variable, for example).
The forward solution of \eqref{PDE} is then computed and the proposal is
accepted or rejected based on some functional of the proposed solution.
The vast majority of proposals are rejected,
making the solution of the discrete forward problems the major
computational bottleneck in the sampling
procedure. At each step in the MCMC method, a discretized version of
\eqref{PDE} must be solved. This leads to a sequence
of linear systems with changing matrix coefficients:

\begin{equation}\label{discrete problem}
A(\omega) x = b.
\end{equation}

If the correlation lengths in the stochastic process defining $a$ are
small compared to the size of the domain,
the solution of even a single realization of \eqref{discrete problem} is
nontrivial and may require adaptivity.
Our goal is: assuming that we have constructed an efficient multilevel
solver for (2) based on some realization of $A(\omega)$, we would like to
efficiently construct a multilevel solver for \eqref{discrete problem}
based on some new realization, $A({\widehat \omega})$.
The nature of MCMC makes this possible because consecutive realizations
of the background $a$ are often similar.

In this talk, we will describe one such adaptive two--level solver in the
setting of smoothed aggregation (or SA) algebraic multigrid (AMG) which
is outlined as follows.
Based on a given hierarchy for $A(\omega)$,
we build a respective two-level solver for the matrix $A({\widehat
\omega})$ corresponding to the
new parameter ${\widehat \omega}$. We test the old hierarchy by iterating
on the homogeneous problem $A({\widehat \omega}) e = 0$. If convergence
is acceptable we retain the old hierarchy. If convergence stalls, then
after several iterations on the homogeneous problem an algebraically
smooth error mode, $x_{\text{bad}}$, has been exposed. This mode is added
to the coarse space to improve performance.
This is achieved by first localizing the problem to every aggregate
exploiting an ``element-free'' approach.
For the localized quadratic forms, we solve generalized eigenvalue
problem. The cost of this computation is greatly reduced by searching
only in the subspace spanned by $x_\text{bad}$ and the basis functions
from the original hierarchy. The choice of this subspace is motivated by
the fact that $A({\omega})$ and $A({\widehat \omega})$, obtained from
MCMC, are in some sense nearby, hence some of the old basis functions may
be adequate whereas the missing part is provided by $x_{\text{bad}}$. The
new (unsmoothed) interpolant is then built based on the first few of the
low-lying modes. The adaptive procedure is repeated as needed.


\end{document}
